\chapter{阅读路线与可复现实验}\label{app:exercises}
\section{三种阅读节奏}
\textbf{先建立直觉。}读第一章和第二章的双峰例子、条件投影框、标准高斯的 SDE/ODE 对照，再读第三章的四轴分类，最后进入第四章的网络和应用。这条路线先获得各对象的定位，随后补推导。

\textbf{以数学为主。}按第二章顺序阅读，做五道章末练习，再补附录 A--C。推导中遇到障碍时，先检查条件和时间方向；不要先试图背最终公式。最后用第三章检查自己能否把公式变成可运行算法。

\textbf{以系统为主。}第一章后进入第四章的 U-Net、隐空间、DiT 与多模态信息流，再回第三章辨认实际系统的路径、目标、求解器和网络，最后回第二章核验公式。仅看到一个网络结构图，不足以确定模型使用哪一种生成动力学。

\section{六个无需大模型的实验}
这些实验是进一步学习建议，并不是本书声称已经训练完成的模型结果。它们可以用一维或二维数据完成，让你把数学关系看成数值事实。

\begin{enumerate}
\item \textbf{DDPM 边缘核验。}对固定 $x_0$，一组样本逐步加噪，另一组直接用式\eqref{eq:forward-marginal}生成。比较单时刻均值、方差与直方图，再比较跨时刻协方差。你会看见边缘一致而联合构造可不同。
\item \textbf{双峰 score。}使用第二章解析混合分布，在不同噪声尺度上绘制密度和 score。找出 score 零点，区分密度极大值与极小值。随后用精确场生成，避免把网络训练误差和动力学原理混在一起。
\item \textbf{SDE/ODE 同边缘。}从标准高斯出发模拟 VP SDE，并比较零速度概率流 ODE。画逐时刻直方图与单条轨迹：前者接近相同，后者明显不同。
\item \textbf{高斯 Flow Matching。}用练习 2.4 的解析速度求解 ODE，比较 Euler、Heun 与解析解。对步数作误差曲线，检查阶数成立的区间。
\item \textbf{参数化权重。}在相同批次与日程上计算噪声损失、原图损失及转换后的权重。逐时刻核验代数关系，观察未经权重修正的损失如何关注不同尺度。
\item \textbf{条件投影。}构造两组在中间位置相交的二维插值，统计相同小邻域里的速度标签和均值。把个别条件直线与边缘速度积分轨迹一起画出，检查它们为何不同。
\end{enumerate}

\section{阅读论文时用同一张问题清单}
每读一篇新论文，先记录：数据表示是什么；端点与概率路径是什么；网络预测哪个目标；损失权重怎样设计；条件怎样注入；采样如何积分；是否另有蒸馏阶段；作者的实验控制了哪些因素。随后再读结果表。

当两篇论文给出不同的结论时，先检查数据、表示、训练预算、引导、评估协议与推理成本是否相同。不要把更大的模型或更多数据带来的收益全归因于一个新公式。第四章的精选阅读路线提供了入口；参考文献中的链接指向论文或官方披露来源。

\section{使用与编译这份书稿}
本书主文件是 \path{main.tex}，正文与附录分别存于独立文件，模板样式在 \path{book-style.tex}。推荐 XeLaTeX 与 Biber。完整顺序为 XeLaTeX、Biber、XeLaTeX、XeLaTeX；最后两遍建立引用和目录。Windows 本地可运行附带的 \path{scripts/build.ps1}。

项目保留 ElegantBook 原始类和许可证，私有依赖只是补齐本地环境所缺的公开组件。在完整 TeX Live 或 Overleaf 环境中通常无需这些额外依赖。具体编译与质量检查记录保存在项目说明中。修改章节后应重新构建整本书，检查未定义引用、长公式和页面溢出，避免只更新源码而仍阅读旧 PDF。
